% Cover letter -- Foundations of Physics (source of the PDF; content mirrors
% COVER_WEN_FOUNDATIONS.md). Compile: pdflatex COVER_WEN_FOUNDATIONS.tex
\documentclass[11pt]{article}
\usepackage[a4paper,margin=1.1in]{geometry}
\usepackage[T1]{fontenc}
\usepackage[utf8]{inputenc}
\usepackage{amsmath,amssymb}
\usepackage[hidelinks]{hyperref}
\usepackage{parskip}
\pagestyle{empty}

\begin{document}

\begin{flushright}
Miqueias Alves Mendes\\
Independent Researcher\\
Ibiapina, Cear\'a, Brazil\\
\href{mailto:mikalvesm@gmail.com}{mikalvesm@gmail.com}\\[4pt]
\today
\end{flushright}

\vspace{6pt}
\noindent To the Editors of \emph{Foundations of Physics},

\vspace{6pt}
I submit for your consideration the manuscript \textbf{``A classification
theorem for emergent $U(1)$ on causal substrates: string-net admissibility
coincides with broken Lorentz invariance.''}

\textbf{What the paper does.} Two research programmes derive gauge structure
rather than postulate it: in condensed matter, string-net condensation
produces emergent photons as collective modes of loop degrees of freedom; in
quantum gravity, causal set theory replaces spacetime by a Lorentz-invariant
discrete order---and no causal substrate has ever produced an emergent
photon. The paper shows these two facts have a single explanation, and states
it as a classification theorem. Define a substrate to be \emph{string-net
admissible} (SNA) if its Hasse diagram offers the three structures string-net
condensation requires---finite valence, two-dimensional amenable ball growth,
and a percolating plaquette density. Then \textbf{a causal substrate is SNA
if and only if it is non-invariant}: every Lorentz-invariant or
covariant-random class fails the predicate by a proven mechanism (a
Campbell--Mecke divergence on the hyperboloid; the non-existence of locally
finite Hasse diagrams for exchangeable orders; a geometric
dichotomy---hyperbolic or one-dimensional---for covariant sequential growth),
while deterministic crystals and foliated lattices realize it. A unified
decision table, with one measured representative per invariance class and no
counterexample, accompanies the theorem. The exact boundary between the two
emergence worlds is the \textbf{definiteness of the metric signature}:
separation orbits are compact spheres for $O(d)$ and infinite-volume quadrics
for $O(p,q)$, so divergence is \emph{chosen} on the Euclidean side and
\emph{forced} on the Lorentzian side. Wick rotation is precisely the
operation that crosses the boundary, and an order cannot cross it and remain
an order.

\textbf{Why \emph{Foundations of Physics}.} The result is kinematic in the
same sense in which Wigner's classification is kinematic: a statement about
what a class of structures can carry, prior to any dynamics. It answers a
foundational question---can gauge structure emerge from causal order
alone?---with a theorem rather than a survey, it explains a long-standing
negative on both sides of the condensed-matter/quantum-gravity divide, and it
is falsifiable at three named joints stated in the paper. This is the
journal's home territory: conceptual structure of physical theories, argued
at theorem grade and kept accountable to measurement.

\textbf{Why this is a complement, not an attack.} The paper does not dispute
string-net theory; it explains why its mechanism lives entirely on the
definite-signature side of the boundary, and why sixty years of
condensed-matter emergence phenomenology do not transport to
Lorentz-invariant causal substrates. Any future emergent-photon claim is told
exactly which lemma it must break.

\textbf{Relation to other work.} Two companion manuscripts reporting the
measured programme (the emergent Newtonian sector and soliton matter; the
$SU(3)$ colour sector on the same substrate) have been submitted separately
to \emph{Physical Review D}; the present paper is the theoretical explanation
of that map, and each stands alone. The measured inputs are cited with their
public source repositories.

\emph{Note on manuscript class.} The source is prepared in REVTeX; I will
supply a Springer-class version at whatever stage the editorial process
requires it.

The manuscript is original and is not under consideration elsewhere. I
declare no conflicts of interest.

Thank you for considering the manuscript.

\vspace{10pt}
\noindent Sincerely,

\vspace{4pt}
\noindent Miqueias Alves Mendes
\end{document}
